% =============================================================================
% Skripsi — format monograf Universitas Indonesia 2025
% Kompilasi: pdfLaTeX -> Biber -> pdfLaTeX -> pdfLaTeX  (Overleaf: otomatis)
% =============================================================================
\documentclass[12pt,a4paper,oneside]{report}
\usepackage{ui-thesis}
\addbibresource{referensi.bib}   % ekspor otomatis Zotero (Better BibTeX)
\addbibresource{manual.bib}      % entri tulisan tangan, tidak ditimpa Zotero

% ---------------------------------------------------------------------------
% DATA TUGAS AKHIR — isi semua baris di bawah
% Judul: tanpa singkatan (kecuali nama/istilah), bukan kalimat tanya, tanpa
% tanda baca penutup (pedoman 10.8).
% ---------------------------------------------------------------------------
\judul{PERBANDINGAN SINGLE-MODEL DAN LLM COUNCIL PADA AGEN VERIFIKASI KERENTANAN DALAM SISTEM LLM MULTI-AGENT LOKAL}
\judulinggris{COMPARING A SINGLE MODEL AND AN LLM COUNCIL FOR THE VULNERABILITY VERIFICATION AGENT OF A LOCAL MULTI-AGENT LLM SYSTEM}
\nama{NAMA LENGKAP}
\npm{NPM}
\jenis{SKRIPSI}
\gelar{Sarjana Teknik}            % cek gelar yang berlaku di prodi
\fakultas{TEKNIK}                 % cek
\departemen{TEKNIK ELEKTRO}       % cek
\prodi{TEKNIK KOMPUTER}           % cek
\kota{DEPOK}
\bulantahun{BULAN TAHUN}          % bulan dan tahun pengesahan, mis. JULI 2027
\tahunakademik{2026/2027}
\pembimbing{Nama Pembimbing, gelar}
\pembimbingdua{}                  % kosongkan jika hanya satu pembimbing
\pengujisatu{Nama Penguji 1, gelar}
\pengujidua{Nama Penguji 2, gelar}
\tanggalsidang{.................}

\begin{document}

% ---------------------------------------------------------------------------
% BAGIAN AWAL (urutan sesuai pedoman 10.1)
% ---------------------------------------------------------------------------
% !TEX root = ../main.tex
% Halaman Sampul (pedoman 10.8): tanpa nomor, satu spasi, tengah, huruf kapital.
\thispagestyle{empty}
\begin{center}
\setstretch{1}\bfseries
\vspace*{0.5cm}
\logoui\\[0.8cm]
UNIVERSITAS INDONESIA\\[3cm]
\thejudul\\[3cm]
\thejenis\\[3.5cm]
\thenama\\[0.3cm]
\thenpm
\vfill
FAKULTAS \thefakultas\\
PROGRAM STUDI \theprodi\\
\thekota\\
\thebulantahun
\end{center}
                % tanpa nomor
\pagenumbering{roman}
\uifrontmatterplain
% !TEX root = ../main.tex
% Halaman Judul (pedoman 10.9): sama dengan sampul + tujuan penyusunan.
% Dihitung sebagai halaman i, tetapi nomornya tidak ditampilkan.
\thispagestyle{empty}
\begin{center}
\setstretch{1}\bfseries
\vspace*{0.5cm}
\logoui\\[0.8cm]
UNIVERSITAS INDONESIA\\[3cm]
\thejudul\\[3cm]
\thejenis\\[0.8cm]
{\normalsize Diajukan sebagai salah satu syarat memperoleh gelar \thegelar}\\[2.5cm]
\thenama\\[0.3cm]
\thenpm
\vfill
FAKULTAS \thefakultas\\
PROGRAM STUDI \theprodi\\
\thekota\\
\thebulantahun
\end{center}
         % i, nomor tidak ditampilkan
% !TEX root = ../main.tex
% Halaman Pernyataan Orisinalitas (pedoman 10.10): spasi ganda, tengah.
% Diunggah sebagai berkas kelengkapan, TIDAK dimasukkan ke repositori UI.
\begin{center}
\bfseries
\MakeUppercase{Halaman Pernyataan Orisinalitas}\\[1.5cm]
\setstretch{2}
Skripsi ini adalah hasil karya saya sendiri,\\
dan semua sumber baik yang dikutip maupun dirujuk\\
telah saya nyatakan dengan benar.
\end{center}
\vspace{3cm}
\begin{center}
\bfseries\setstretch{2}
\begin{tabular}{@{}l@{\hspace{2cm}}l@{}}
Nama & : \thenama \\
NPM  & : \thenpm \\
Tanda Tangan & : \\[1.5cm]
Tanggal & : \thetanggalsidang \\
\end{tabular}
\end{center}
          % ii
% !TEX root = ../main.tex
% Halaman Pernyataan Penggunaan Kecerdasan Artifisial (pedoman 4.1 butir 9, 10.11).
% Pilih SALAH SATU pernyataan dan hapus yang lain. Jika memakai AI, sebutkan
% platform dan tujuannya dengan jujur dan dapat dipertanggungjawabkan.
\begin{center}
\bfseries\setstretch{2}
HALAMAN PERNYATAAN PENGGUNAAN\\
KECERDASAN ARTIFISIAL
\end{center}
\vspace{1cm}
{\bfseries\setstretch{2}\sloppy
% --- Pilihan 1: tidak memakai AI ---
% \begin{center}
% Saya menyatakan bahwa dalam penyusunan Tugas Akhir ini, saya tidak
% menggunakan teknologi Kecerdasan Buatan dalam bentuk apa pun.
% \end{center}
% --- Pilihan 2: memakai AI ---
\noindent Saya menyatakan bahwa dalam penyusunan Tugas Akhir ini, saya
menggunakan teknologi Kecerdasan Buatan dengan platform \underline{\hspace{3cm}}
untuk tujuan \underline{\hspace{3cm}} (misal pemeriksaan dan perbaikan tata bahasa)
dan \underline{\hspace{3cm}} dan saya bertanggung jawab penuh atas keakuratan,
orisinalitas, dan integritas seluruh isi Tugas Akhir ini.\par}
\vspace{1.5cm}
\begin{center}
\bfseries\setstretch{2}
\begin{tabular}{@{}l@{\hspace{2cm}}l@{}}
Nama & : \thenama \\
NPM  & : \thenpm \\
Tanda Tangan & : \\
Tanggal & : \thetanggalsidang \\
\end{tabular}
\end{center}
         % iii
% !TEX root = ../main.tex
% Halaman Pengesahan (pedoman 10.12): satu spasi.
% Diunggah sebagai berkas kelengkapan, TIDAK dimasukkan ke repositori UI.
\begin{center}\bfseries HALAMAN PENGESAHAN\end{center}
\vspace{0.5cm}
{\setstretch{1}
\noindent Skripsi ini diajukan oleh:\\[0.3cm]
\begin{tabular}{@{}p{4cm}@{\,:\,}>{\raggedright\arraybackslash}p{9.5cm}@{}}
Nama & \thenama \\
NPM & \thenpm \\
Program Studi & \theprodi \\
Judul Skripsi & \thejudul \\
\end{tabular}

\vspace{1cm}
\noindent\textbf{Telah berhasil dipertahankan di hadapan Dewan Penguji dan diterima
sebagai bagian persyaratan yang diperlukan untuk memperoleh gelar \thegelar{} pada
Program Studi \theprodi, Fakultas \thefakultas, Universitas Indonesia.}

\vspace{1.5cm}
\noindent\textbf{DEWAN PENGUJI}\\[0.6cm]
\begin{tabular}{@{}l@{\ :\ }p{7.5cm}l@{}}
Pembimbing & \thepembimbing & (\hspace{3cm}) \\[1cm]
\barispembimbingdua
Penguji & \thepengujisatu & (\hspace{3cm}) \\[1cm]
Penguji & \thepengujidua & (\hspace{3cm}) \\
\end{tabular}

\vfill
\noindent Ditetapkan di : \thekota\\
Tanggal\hspace{1.2cm}: \thetanggalsidang
\par}
            % iv
% !TEX root = ../main.tex
% Surat Pernyataan bebas plagiat (ada di Template Tugas Akhir UI 2025).
% Hapus \include-nya di main.tex jika fakultas tidak memakai halaman ini.
\begin{center}\bfseries SURAT PERNYATAAN\end{center}
\vspace{0.5cm}
\noindent Yang bertanda tangan di bawah ini, saya:\\[0.3cm]
\begin{tabular}{@{}p{4cm}@{\,:\,}l@{}}
Nama & \thenama \\
NPM & \thenpm \\
Program Studi & \theprodi \\
Tahun Akademik & \thetahunakademik \\
\end{tabular}

\noindent Menyatakan bahwa saya tidak melakukan kegiatan plagiat dalam penulisan
skripsi saya yang berjudul:
\begin{center}
\setstretch{1.2}\bfseries ``\thejudul''
\end{center}
\noindent Apabila suatu saat nanti terbukti saya melakukan plagiat, maka saya akan
menerima sanksi yang telah ditetapkan.

Demikian pernyataan ini saya buat dengan sebenarnya.

\vspace{1cm}
\noindent \thekota, \thetanggalsidang\\[2cm]
\noindent (\thenama)
      % ada di Template UI 2025; hapus jika fakultas tidak memakai
% !TEX root = ../main.tex
% Kata Pengantar / Ucapan Terima Kasih (pedoman 10.13):
% spasi 1,5; urutan terima kasih: pihak internal UI, pihak luar, lalu keluarga/teman.
\chapter*{KATA PENGANTAR}
\phantomsection\addcontentsline{toc}{chapter}{KATA PENGANTAR}

Puji syukur saya panjatkan kepada ... . Penulisan skripsi ini dilakukan dalam
rangka memenuhi salah satu syarat untuk mencapai gelar \thegelar{} Program Studi
\theprodi{} pada Fakultas \thefakultas{} Universitas Indonesia. Saya menyadari bahwa,
tanpa bantuan dan bimbingan dari berbagai pihak, sangatlah sulit bagi saya untuk
menyelesaikan skripsi ini. Oleh karena itu, saya mengucapkan terima kasih kepada:
\begin{enumerate}
  \item \thepembimbing, selaku dosen pembimbing yang telah ...;
  \item ... (pihak internal UI lainnya);
  \item ... (pihak luar UI);
  \item orang tua dan keluarga saya yang telah ...; dan
  \item sahabat yang telah ... .
\end{enumerate}

Akhir kata, saya berharap ... . Semoga skripsi ini membawa manfaat bagi
pengembangan ilmu.

\vspace{1cm}
\begin{flushright}
\thekota, \thetanggalsidang\\[1.5cm]
\thenama
\end{flushright}

% !TEX root = ../main.tex
% Halaman Pernyataan Persetujuan Publikasi (pedoman 10.14): spasi 1,5; judul satu spasi.
% Diunggah sebagai berkas kelengkapan, TIDAK dimasukkan ke repositori UI.
\begin{center}
\setstretch{1}\bfseries
HALAMAN PERNYATAAN PERSETUJUAN PUBLIKASI\\
TUGAS AKHIR UNTUK KEPENTINGAN AKADEMIS
\end{center}
\noindent\rule{\textwidth}{1pt}
\vspace{0.3cm}

\noindent Sebagai sivitas akademik Universitas Indonesia, saya yang bertanda tangan di
bawah ini:\\[0.2cm]
\begin{tabular}{@{}p{3.5cm}@{\,:\,}l@{}}
Nama & \thenama \\
NPM & \thenpm \\
Program Studi & \theprodi \\
Departemen & \thedepartemen \\
Fakultas & \thefakultas \\
Jenis karya & Skripsi \\
\end{tabular}

\noindent demi pengembangan ilmu pengetahuan, menyetujui untuk memberikan kepada
Universitas Indonesia \textbf{Hak Bebas Royalti Noneksklusif (\textit{Non-exclusive
Royalty-Free Right})} atas karya ilmiah saya yang berjudul:
\begin{center}\setstretch{1.2}\thejudul\end{center}
\noindent beserta perangkat yang ada (jika diperlukan). Dengan Hak Bebas Royalti
Noneksklusif ini Universitas Indonesia berhak menyimpan, mengalihmedia/formatkan,
mengelola dalam bentuk pangkalan data (\textit{database}), merawat, dan
memublikasikan tugas akhir saya selama tetap mencantumkan nama saya sebagai
penulis/pencipta dan sebagai pemilik Hak Cipta.

Demikian pernyataan ini saya buat dengan sebenarnya.

\vspace{0.5cm}
\begin{center}
\setstretch{1}
Dibuat di : \thekota\\
Pada tanggal : \thetanggalsidang\\
Yang menyatakan\\[2cm]
(\thenama)
\end{center}


\uifrontmatter                       % footer "Universitas Indonesia" mulai Abstrak
% !TEX root = ../main.tex
% Abstrak (pedoman 10.17, format monograf):
%   - satu paragraf, maksimal 500 kata, satu spasi;
%   - dua bahasa, sebisa mungkin dalam satu halaman;
%   - nama (tanpa NPM), program studi, dan judul di atas abstrak;
%   - kata kunci di bawah, dalam bahasa masing-masing;
%   - semua istilah asing (kecuali nama) dicetak miring.
\chapter*{ABSTRAK}
\phantomsection\addcontentsline{toc}{chapter}{ABSTRAK}
\vspace{-1\baselineskip}
{\setstretch{1}
\noindent\begin{tabular}{@{}p{3cm}@{\,:\,}>{\raggedright\arraybackslash}p{10.5cm}@{}}
Nama & \thenama \\
Program Studi & \theprodi \\
Judul & \thejudul \\
Pembimbing & \thepembimbing \\
\end{tabular}

\vspace{0.4cm}
\noindent Tulis abstrak dalam satu paragraf: masalah dan celah penelitian, tujuan,
metode (konfigurasi \textit{single-model}, \textit{self-consistency}, dan
\textit{LLM council} pada Agen Verifikasi), data, hasil utama dengan angka, dan
kesimpulan. Maksimal 500 kata.

\vspace{0.3cm}
\noindent\textbf{Kata kunci:} kata kunci 1, kata kunci 2, kata kunci 3

\vspace{1cm}
\begin{center}\bfseries ABSTRACT\end{center}
\noindent\begin{tabular}{@{}p{3cm}@{\,:\,}>{\raggedright\arraybackslash}p{10.5cm}@{}}
Name & \thenama \\
Study Program & \theprodi \\
Title & \thejudulinggris \\
Supervisor & \thepembimbing \\
\end{tabular}

\vspace{0.4cm}
\noindent\textit{Write the English abstract here, as a single paragraph of at most
500 words, matching the Indonesian abstract.}

\vspace{0.3cm}
\noindent\textbf{\textit{Keywords:}} \textit{keyword 1, keyword 2, keyword 3}
\par}


\begingroup
\setstretch{1}
\tableofcontents
\clearpage
\phantomsection\addcontentsline{toc}{chapter}{DAFTAR GAMBAR}
\listoffigures
\clearpage
\phantomsection\addcontentsline{toc}{chapter}{DAFTAR TABEL}
\listoftables
\endgroup

% ---------------------------------------------------------------------------
% ISI
% ---------------------------------------------------------------------------
\uimainmatter
% !TEX root = ../main.tex
\chapter{PENDAHULUAN}
\label{bab:pendahuluan}

\section{Latar Belakang}
\label{sec:latar-belakang}
% Empat langkah (pedoman 9.1): (1) situasi, (2) mengapa penting, (3) celah
% penelitian, (4) apa yang dilakukan skripsi ini. Setiap klaim faktual diberi
% sitasi. Akhiri dengan celah, bukan dengan latar belakang umum.

% (1) Situasi: AI mempercepat penemuan kerentanan, jumlah CVE melonjak.
Kecerdasan buatan (\asing{artificial intelligence}, AI) kini tidak hanya
dipakai untuk menulis perangkat lunak, tetapi juga untuk mencari celah keamanan
di dalamnya. Agen berbasis \asing{large language model} (LLM) mampu menelusuri
kode sumber, menyusun hipotesis kerentanan, lalu membuktikannya dengan
\asing{proof of concept}~\cite{addisLLMBasedIntelligentAgents2026}. Kemampuan
ini diuji dalam skala besar pada \asing{AI Cyber Challenge} (AIxCC) yang
diselenggarakan DARPA. Pada babak final Agustus 2025, tujuh sistem otonom
menganalisis 53 proyek perangkat lunak sumber terbuka selama sekitar 143 jam
tanpa campur tangan manusia, dan seluruh tim menemukan setidaknya satu
kerentanan \asing{zero-day} yang sebelumnya tidak diketahui~\cite{zhang2026aixcc}.
Di luar kompetisi, pengembang model juga mulai memindai perangkat lunak sumber
terbuka secara massal. Anthropic melaporkan bahwa pemindainya telah menemukan
lebih dari 29.000 kandidat kerentanan, tetapi baru sekitar 6.000 di antaranya
yang sempat ditinjau secara manual karena kapasitas manusia untuk memvalidasi
temuan menjadi hambatan utama~\cite{anthropic2026ossscanner}. Pada tolok ukur
CyberGym, LLM yang pada awal 2025 menemukan kurang dari 20\% kerentanan kini
menemukan lebih dari 85\%~\cite{anthropic2026ossscanner}.

Percepatan penemuan ini tampak pada jumlah \asing{Common Vulnerabilities and
Exposures} (CVE) yang dipublikasikan. Pada 2024 tercatat sekitar 40.000 CVE,
pada 2025 jumlahnya mencapai rekor 48.185 CVE, dan pada 2026 diproyeksikan
sekitar 66.000 CVE~\cite{hand2026vulnpocalypse}. Penemuan kerentanan otomatis
oleh model AI disebut sebagai salah satu pendorong lonjakan tersebut, di
samping bertambahnya kode yang ditulis oleh AI itu
sendiri~\cite{hand2026vulnpocalypse}. Lembaga yang mengelola basis data
kerentanan pun kewalahan. NIST melaporkan kenaikan pengajuan sebesar 263\%
sejak 2020 dan kini memprioritaskan analisis berdasarkan tingkat risiko, bukan
menganalisis setiap CVE; pada 2025 hanya sekitar 28\% CVE baru yang dianalisis
penuh oleh \asing{National Vulnerability Database}
(NVD)~\cite{hand2026vulnpocalypse}. Pada saat yang sama, jeda antara
publikasi kerentanan dan eksploitasinya semakin pendek, sehingga CISA
mengganti tenggat perbaikan tetap 14 hari dengan model berbasis risiko yang
dapat menuntut tindakan dalam tiga hari~\cite{hand2026vulnpocalypse}.

Kemampuan AI yang sama juga sudah dipakai oleh penyerang. ARTEX, sebuah
\asing{framework} \asing{penetration testing} sumber terbuka yang menjalankan
beberapa agen LLM secara otonom, telah disalahgunakan oleh berbagai kelompok
peretas sejak Agustus 2026 untuk menyerang organisasi di sektor pemerintahan,
telekomunikasi, keuangan, dan manufaktur~\cite{gtig2026artex}. Operator cukup
memberikan tujuan dalam bahasa alami, lalu agen-agen ARTEX menjalankan
pengintaian, eksploitasi, eskalasi hak akses, dan perpindahan ke sistem lain
secara mandiri. Dalam salah satu kampanye, pelaku bergerak dari akun pertama
yang dibobol hingga akses sistem yang dalam dalam waktu sekitar empat jam, dan
pada satu kasus hanya 15 menit~\cite{gtig2026artex}. FBI mencatat bahwa
sebagian intrusi mencapai akses istimewa dalam waktu kurang dari tiga jam sejak
penargetan awal, dan bahwa ARTEX memungkinkan pelaku dengan keahlian teknis
minim menjalankan operasi yang kompleks terhadap banyak target
sekaligus~\cite{fbi2026artex}. ARTEX juga sangat efektif menyerang sistem
pinggiran yang kurang mendapat perhatian keamanan, seperti portal mitra dan
pemasok, \asing{backend} aplikasi seluler, dan layanan lama, serta memanfaatkan
vendor pihak ketiga yang lebih lemah sebagai pintu masuk ke jaringan perusahaan
yang lebih besar~\cite{fbi2026artex}. Hingga awal November 2026, lebih dari 1.600
konsol ARTEX ditemukan terbuka di internet~\cite{fbi2026artex}. CERT-EU menilai bahwa
ARTEX mampu menguji sejumlah besar salah konfigurasi sederhana yang tidak
sempat diperiksa oleh penguji manusia dan merangkainya menjadi jalur serangan,
sehingga menyarankan agar setiap temuan, termasuk yang tidak parah, tetap
dimasukkan ke dalam siklus perbaikan~\cite{certeu2026artex}. Dengan demikian,
kerentanan tidak hanya ditemukan semakin cepat, tetapi juga dieksploitasi
semakin cepat.

% (2) Mengapa penting: UMKM tanpa tim keamanan siber kewalahan.
Bagi organisasi besar yang memiliki pusat operasi keamanan, lonjakan ini
menambah beban kerja. Bagi usaha kecil dan menengah (UKM), lonjakan ini dapat
melampaui kemampuan mereka sama sekali. UKM umumnya tidak memiliki tim keamanan
siber khusus, anggaran yang memadai, maupun pengetahuan untuk menilai sendiri
kerentanan pada sistem yang mereka
jalankan~\cite{chaudharyQuestResearchKnowledge2023,michalekCyberIncidentPrevention2026}.
Padahal, dari puluhan ribu CVE setiap tahun, hanya sebagian kecil yang benar-benar
relevan dan dapat dieksploitasi pada sistem tertentu: data EPSS menunjukkan bahwa
hanya sekitar 5\% CVE yang pernah teramati dieksploitasi~\cite{hand2026vulnpocalypse}.
Pekerjaan memilah CVE yang relevan, memastikan apakah kerentanan itu benar-benar
ada pada sistem, lalu memprioritaskan perbaikannya merupakan inti manajemen
kerentanan~\cite{syedCybersecurityVulnerabilityManagement2020,shimizuVulnerabilityManagementChaining2026}.
Tanpa tim khusus, pekerjaan ini praktis tidak dilakukan, sehingga kerentanan
yang sebenarnya dapat diperbaiki dibiarkan terbuka.

Sistem multi-agen berbasis LLM menawarkan jalan keluar dengan mengotomatiskan
alur manajemen kerentanan. Wijaya~\cite{wijaya2026} merancang sistem semacam ini
menggunakan n8n sebagai orkestrator. Sistem tersebut terdiri atas Agen Direktur,
Agen Pencari, Agen Peneliti, Agen Verifikasi, dan Agen Pengamat Teknis yang
bersama-sama mengidentifikasi layanan pada server target, mencari CVE yang
relevan, menyusun spesifikasi teknis, lalu membuat dan menjalankan kode
verifikasi untuk memastikan keberadaan kerentanan. Namun, sistem tersebut
memiliki tiga keterbatasan yang saling berkaitan~\cite{wijaya2026}. Pertama,
sebagian besar agen bergantung pada LLM berbasis \asing{cloud}, sehingga
informasi sensitif organisasi, seperti layanan yang berjalan, konfigurasi sistem,
dan kerentanan yang ditemukan, dikirim ke pihak ketiga. Risiko kebocoran dan
penyalahgunaan data semacam ini merupakan persoalan keamanan tersendiri dalam
operasional LLM di \asing{cloud}~\cite{aliDoubleedgedSwordLLM2026}. Kedua, LLM
publik cenderung menolak membuat kode yang menyerupai eksploit, padahal kode
semacam itu diperlukan untuk verifikasi. Ketiga, Agen Verifikasi yang dijalankan
secara lokal dengan model kecil di CPU memakan sekitar 57\% waktu eksekusi dan
tidak berhasil memverifikasi satu kerentanan pun. Dengan kata lain, Agen
Verifikasi merupakan titik terlemah sistem tersebut.

Dua keterbatasan pertama mendorong agar seluruh sistem dijalankan secara
\asing{self-hosted} dengan model \asing{open-weight}. Pendekatan ini tidak
otomatis menjamin privasi~\cite{alzoubiAssessingPrivacyRisks2026} dan menuntut
infrastruktur yang sering kali tidak dimiliki UKM~\cite{Islam_2026}, tetapi
setidaknya menjaga data kerentanan tetap berada di dalam organisasi. Keterbatasan
ketiga menunjukkan bahwa \asing{self-hosting} membawa kendala sumber daya. Dari
sisi \asing{serving}, vLLM dengan \asing{PagedAttention} meningkatkan
\asing{throughput} 2--4 kali dibandingkan sistem
sebelumnya~\cite{kwonEfficientMemoryManagement} dan terbukti jauh lebih stabil
daripada Ollama pada beban permintaan bersamaan~\cite{ayBenchmarkingOllamaVLLM2026},
sehingga layak menggantikan Ollama yang dipakai pada sistem
dasar~\cite{wijaya2026}. Dari sisi kualitas, model lokal yang muat pada satu GPU
umumnya lebih lemah daripada model proprietary
terbesar~\cite{koterbaLeveragingLargeLanguage2026}, sehingga diperlukan cara lain
untuk meningkatkan ketepatan verifikasi.

% (3) Celah penelitian.
Salah satu cara tersebut adalah menggabungkan keluaran beberapa LLM, yang dalam
penelitian ini disebut \asing{LLM council}. Gagasannya, beberapa model yang
berbeda mengerjakan tugas yang sama, lalu hasilnya digabungkan melalui
\asing{voting} atau dipilih oleh satu model penilai (\asing{chairman}). Survei
\asing{ensemble} LLM menunjukkan bahwa penggabungan semacam ini dapat
meningkatkan kualitas pembangkitan teks dan kode~\cite{ashigaEnsembleLearningLarge2026},
dan sistem multi-agen yang tersusun dari model kecil mampu menyaingi model
besar pada pembangkitan kode~\cite{zhangPowerSmallLLMs2026}. Pada domain
kerentanan, sistem multi-agen LLM juga telah dipakai untuk analisis kerentanan
perangkat IoT~\cite{yangLIVAMultiAgentLLMAssisted}. Namun, penggabungan beberapa
model tidak selalu menguntungkan. Halusinasi satu agen dapat menular ke agen
lain~\cite{zhangBewareWoozleEffect2026}, dan setiap anggota tambahan
melipatgandakan token, latensi, memori GPU, dan energi yang dibutuhkan. Bagi UKM
yang hanya mampu menyediakan satu server GPU, biaya tambahan ini menentukan
apakah sebuah pendekatan layak dipakai atau tidak.

Penelusuran literatur pada jurnal IEEE dan ScienceDirect periode 2023--2026
tidak menemukan penelitian yang membandingkan model tunggal dengan \asing{LLM
council} untuk verifikasi kerentanan pada model lokal, apalagi yang mengukur
pertukaran (\asing{trade-off}) antara peningkatan akurasi dan tambahan
penggunaan sumber daya. Selain itu, verifikasi pada sistem
dasar~\cite{wijaya2026} menghasilkan kode yang kemudian dijalankan pada server
target, bukan jawaban teks pendek. Akibatnya, teknik penggabungan yang lazim,
seperti \asing{majority vote} atas jawaban, tidak dapat langsung diterapkan,
dan belum diketahui strategi penggabungan mana yang paling sesuai untuk
keluaran berupa kode.

% (4) Apa yang dilakukan penelitian ini.
Penelitian ini menjalankan sistem multi-agen manajemen kerentanan
Wijaya~\cite{wijaya2026} secara lokal dengan vLLM sebagai \asing{serving engine}
dan n8n sebagai orkestrator, lalu membandingkan beberapa konfigurasi Agen
Verifikasi: model tunggal, \asing{self-consistency} dengan model tunggal yang
sama sebagai kontrol beban komputasi, dan \asing{LLM council} dengan dua
strategi penggabungan, yaitu \asing{voting} atas hasil eksekusi kode dan
pemilihan kode oleh \asing{chairman}. Setiap konfigurasi diukur dari sisi
akurasi terhadap data uji berlabel dan dari sisi penggunaan sumber daya, yaitu
token, latensi, VRAM, dan energi GPU. Hasilnya diharapkan menjawab apakah
peningkatan akurasi yang diberikan \asing{LLM council}, jika ada, sepadan dengan
sumber daya tambahan yang dibutuhkannya pada perangkat keras yang realistis bagi
UKM.

\section{Rumusan Masalah}
\label{sec:rumusan-masalah}
% Satu masalah utama dan 2--3 submasalah (pedoman 9.1 butir 2).
Masalah utama penelitian ini adalah belum diketahuinya apakah \asing{LLM council}
memberikan peningkatan akurasi verifikasi kerentanan yang sepadan dengan tambahan
penggunaan sumber daya dibandingkan model tunggal pada sistem multi-agen berbasis
LLM lokal. Masalah tersebut dirinci menjadi tiga pertanyaan berikut.
\begin{enumerate}
  \item Apakah \asing{LLM council} menghasilkan verifikasi kerentanan yang lebih
        akurat (\asing{precision}, \asing{recall}, F1, dan \asing{false-positive
        rate}) dibandingkan model tunggal terbaik dan \asing{self-consistency}
        pada perangkat keras lokal yang sama?
  \item Berapa tambahan penggunaan sumber daya (token, latensi, VRAM, dan energi)
        yang dibutuhkan \asing{LLM council}, dan bagaimana pertukaran antara
        akurasi dan sumber daya tersebut dibandingkan model tunggal?
  \item Strategi penggabungan mana yang lebih baik untuk keluaran berupa kode
        verifikasi: \asing{voting} atas hasil eksekusi atau pemilihan kode oleh
        \asing{chairman} sebelum eksekusi?
\end{enumerate}

\section{Hipotesis Penelitian}
\label{sec:hipotesis}
% Satu hipotesis per rumusan masalah, dapat diuji oleh eksperimen.
\begin{enumerate}
  \item \asing{LLM council} dengan model yang beragam menghasilkan F1 yang lebih
        tinggi dan \asing{false-positive rate} yang lebih rendah dibandingkan
        model tunggal terbaik maupun \asing{self-consistency} dengan jumlah
        sampel yang sama.
  \item \asing{LLM council} membutuhkan token, latensi, VRAM, dan energi per kasus
        yang lebih besar dibandingkan model tunggal, tetapi energi per verifikasi
        yang benar tidak meningkat sebanding dengan jumlah anggota
        \asing{council}.
  \item \asing{Voting} atas hasil eksekusi menghasilkan akurasi yang lebih tinggi
        dibandingkan pemilihan kode oleh \asing{chairman}, dengan konsekuensi
        jumlah eksekusi pada server target dan penggunaan sumber daya yang lebih
        besar.
\end{enumerate}

\section{Tujuan Penelitian}
\label{sec:tujuan}
% Satu tujuan per rumusan masalah, dengan urutan yang sama (pedoman 9.1 butir 3).
\begin{enumerate}
  \item Membandingkan akurasi verifikasi kerentanan antara model tunggal,
        \asing{self-consistency}, dan \asing{LLM council} pada Agen Verifikasi
        yang berjalan secara lokal.
  \item Mengukur penggunaan sumber daya setiap konfigurasi dan menganalisis
        pertukaran antara akurasi dan sumber daya tersebut.
  \item Membandingkan efektivitas dan biaya strategi \asing{voting} atas hasil
        eksekusi dengan strategi pemilihan kode oleh \asing{chairman}.
\end{enumerate}

\section{Manfaat Penelitian}
\label{sec:manfaat}
% Pedoman 9.1 butir 4: manfaat akademis dan manfaat praktis.
\subsection{Manfaat Akademis}
Penelitian ini memberikan bukti empiris tentang efektivitas \asing{LLM council}
pada tugas verifikasi kerentanan yang keluarannya berupa kode yang dapat
dijalankan, bukan jawaban teks. Penggunaan \asing{self-consistency} sebagai
kontrol memisahkan pengaruh keberagaman model dari pengaruh bertambahnya
jumlah sampel. Penelitian ini juga menyajikan pengukuran sumber daya, termasuk
energi GPU, yang dapat dipakai sebagai acuan bagi penelitian lain tentang
penggabungan beberapa LLM pada perangkat keras terbatas.

\subsection{Manfaat Praktis}
Hasil penelitian ini dapat menjadi pertimbangan bagi UKM atau organisasi tanpa
tim keamanan siber khusus yang ingin menjalankan sistem manajemen kerentanan
berbasis AI di infrastruktur sendiri, terutama dalam memilih antara satu model
atau beberapa model dengan anggaran perangkat keras yang terbatas. Penelitian
ini juga menghasilkan versi sistem Wijaya~\cite{wijaya2026} yang seluruhnya
berjalan secara lokal dengan vLLM, sehingga data kerentanan organisasi tidak
perlu dikirim ke layanan pihak ketiga.

\section{Ruang Lingkup Penelitian}
\label{sec:ruang-lingkup}
% Batasan, asumsi, dan keterbatasan (pedoman 9.1 butir 5).
\begin{enumerate}
  \item Perbandingan hanya dilakukan pada Agen Verifikasi. Agen Direktur, Agen
        Pencari, Agen Peneliti, dan Agen Pengamat Teknis dijalankan secara lokal
        tetapi tidak dioptimasi atau dievaluasi secara mendalam.
  \item Orkestrasi memakai n8n seperti pada sistem dasar~\cite{wijaya2026}, dan
        \asing{serving} model memakai vLLM dengan model \asing{open-weight} yang
        muat pada satu GPU kelas konsumen atau \asing{workstation}.
  \item Arsitektur asli sistem dasar dijalankan kembali sebagai kontrol. Energi
        bagian sistem dasar yang berjalan di \asing{cloud} tidak dapat diukur,
        sehingga hanya energi komponen lokal yang dilaporkan untuk konfigurasi
        tersebut.
  \item Pengujian dilakukan pada lingkungan laboratorium virtual dengan server
        target yang menjalankan perangkat lunak dengan kerentanan yang telah
        diketahui, bukan pada sistem produksi. Kebenaran hasil verifikasi dinilai
        terhadap data uji berlabel yang berisi kasus rentan dan tidak rentan.
  \item Penggunaan energi diukur pada GPU menggunakan pustaka Zeus, sedangkan
        energi CPU dan seluruh sistem hanya dilaporkan sebagai pelengkap.
  \item \asing{LLM council} dibatasi pada tiga model dari keluarga yang berbeda
        dengan ukuran yang sebanding.
\end{enumerate}

\section{Sistematika Penulisan}
\label{sec:sistematika}
% Opsional untuk format monograf; satu paragraf singkat per bab.
\textbf{Bab~\ref{bab:pendahuluan} Pendahuluan} berisi latar belakang, rumusan
masalah, hipotesis, tujuan, manfaat, ruang lingkup penelitian, dan sistematika
penulisan.

\textbf{Bab~\ref{bab:tinjauan} Tinjauan Pustaka} berisi teori tentang manajemen
kerentanan dan CVE, LLM dan kuantisasi, \asing{serving} LLM dengan vLLM, sistem
multi-agen, teknik penggabungan keluaran beberapa LLM, serta metrik evaluasi,
dilanjutkan dengan penelitian terkait dan analisis kesenjangan penelitian.

\textbf{Bab~\ref{bab:metodologi} Metodologi Penelitian} berisi arsitektur sistem
dan lingkungan pengujian, konfigurasi Agen Verifikasi yang dibandingkan,
penyusunan data uji berlabel, serta prosedur pengukuran akurasi dan penggunaan
sumber daya.

\textbf{Bab~\ref{bab:hasil} Hasil dan Pembahasan} berisi hasil pengujian setiap
konfigurasi, analisis pertukaran antara akurasi dan sumber daya, serta pembahasan
terhadap hipotesis penelitian.

\textbf{Bab~\ref{bab:kesimpulan} Kesimpulan dan Saran} berisi jawaban atas
rumusan masalah dan saran untuk penelitian selanjutnya.

% !TEX root = ../main.tex
\chapter{TINJAUAN PUSTAKA}
\label{bab:tinjauan}
% Pedoman 9.2: teori, kajian literatur, kesenjangan, kerangka pemikiran,
% definisi operasional. Bab ini harus berakhir pada celah yang diklaim di Bab 1.

\section{Landasan Teori}
\label{sec:teori}
\subsection{Manajemen Kerentanan dan CVE}
...
\subsection{\asing{Large Language Model} dan Kuantisasi}
...
\subsection{\asing{Serving} LLM dengan vLLM}
...
\subsection{Sistem Multi-Agen dan Orkestrasi dengan n8n}
...
\subsection{Penggabungan Jawaban Beberapa LLM}
% Self-consistency, majority/weighted voting, multi-agent debate,
% LLM-as-a-judge/chairman. Bagian teori inti: beri ruang paling banyak.
...
\subsection{Evaluasi Klasifikasi}
% Confusion matrix, precision, recall, F1, Fleiss' kappa, uji McNemar.
...

\section{Penelitian Terkait}
\label{sec:penelitian-terkait}
% Kelompokkan per tema, bukan satu paragraf per paper.
\subsection{\asing{Multi-Agent Debate} dan \asing{LLM Council}}
...
\subsection{\asing{Ensemble} dan \asing{Voting} pada LLM}
...
\subsection{\asing{LLM-as-a-Judge}}
...
\subsection{LLM dan Agen untuk Keamanan Siber}
...
\subsection{Model Kecil dan \asing{Deployment} Lokal}
...

% Tabel perbandingan penelitian terkait (contoh format tabel).
\begin{table}[htbp]
  \caption{Perbandingan penelitian terkait}
  \label{tab:perbandingan}
  \centering\small\setstretch{1}
  \begin{tabularx}{\textwidth}{@{}l X l l X@{}}
    \toprule
    Penelitian & Domain & Jumlah model & Lokal/\asing{cloud} & Hasil utama \\
    \midrule
    ... & ... & ... & ... & ... \\
    Penelitian ini & Verifikasi kerentanan & 1 vs 3 & Lokal & ... \\
    \bottomrule
  \end{tabularx}
\end{table}

\section{Analisis Kesenjangan Penelitian}
\label{sec:kesenjangan}
...

\section{Kerangka Pemikiran}
\label{sec:kerangka}
% Gambar model konseptual: masalah -> intervensi -> hasil yang diukur.
\begin{figure}[htbp]
  \centering
  \fbox{\parbox[c][4cm][c]{0.8\textwidth}{\centering Gambar kerangka pemikiran}}
  \caption{Kerangka pemikiran penelitian}
  \label{fig:kerangka}
\end{figure}

\section{Definisi Operasional}
\label{sec:definisi-operasional}
\begin{table}[htbp]
  \caption{Definisi operasional}
  \label{tab:definisi-operasional}
  \centering\small\setstretch{1}
  \begin{tabularx}{\textwidth}{@{}l X X l@{}}
    \toprule
    Istilah & Definisi & Cara mengukur & Satuan \\
    \midrule
    \asing{True positive} & ... & ... & kasus \\
    \asing{False positive} & ... & ... & kasus \\
    F1 & ... & ... & -- \\
    Latensi p50/p95 & ... & ... & detik \\
    Total token & ... & ... & token \\
    VRAM puncak & ... & ... & GB \\
    Fleiss' kappa & ... & ... & -- \\
    \bottomrule
  \end{tabularx}
\end{table}

% !TEX root = ../main.tex
\chapter{METODOLOGI PENELITIAN}
\label{bab:metodologi}
% Pedoman 9.3: cukup rinci agar eksperimen dapat diulang oleh orang lain.

\section{Jenis dan Desain Penelitian}
\label{sec:desain}
% Kuantitatif, eksperimen komparatif. Sertakan gambar alur penelitian.
...

\section{Rancangan Sistem}
\label{sec:rancangan}
\subsection{Arsitektur Sistem}
% Alur n8n: Direktur -> Pencari -> Peneliti -> Agen Verifikasi -> Pengamat Teknis.
...
\subsection{Lapisan \asing{Serving}}
% vLLM, model dan kuantisasi, cara tiga model berbagi satu GPU.
...
\subsection{Konfigurasi Agen Verifikasi}
\begin{table}[htbp]
  \caption{Konfigurasi Agen Verifikasi yang dibandingkan}
  \label{tab:konfigurasi}
  \centering\setstretch{1}
  \begin{tabularx}{\textwidth}{@{}l X X@{}}
    \toprule
    Kode & Konfigurasi & Fungsi \\
    \midrule
    S1--S3 & Setiap model kandidat sendirian & \asing{Baseline}; mencari \asing{single model} terbaik \\
    SC & \asing{Single model} terbaik, $k$ sampel, \asing{majority vote} & Kontrol: komputasi setara \asing{council} tanpa keberagaman model \\
    C-vote & Tiga model berbeda, \asing{majority vote} & \asing{Council} paling sederhana \\
    C-chair & Tiga model menjawab, satu \asing{chairman} menyintesis & Pola \asing{LLM council} \\
    \bottomrule
  \end{tabularx}
\end{table}
\subsection{\asing{Prompt} dan Format Keluaran}
...

\section{Objek Penelitian dan \asing{Dataset}}
\label{sec:dataset}
...

\section{Lingkungan dan Instrumen Penelitian}
\label{sec:lingkungan}
...

\section{Prosedur Eksperimen}
\label{sec:prosedur}
...

\section{Metode Analisis Data}
\label{sec:analisis}
% Contoh rumus:
\begin{equation}
  F_1 = \frac{2 \cdot \mathit{precision} \cdot \mathit{recall}}{\mathit{precision} + \mathit{recall}}
  \label{eq:f1}
\end{equation}
...

\section{Etika Penelitian}
\label{sec:etika}
...

% !TEX root = ../main.tex
\chapter{HASIL DAN PEMBAHASAN}
\label{bab:hasil}
% Pedoman 9.4: sajikan data sesuai urutan rumusan masalah, objektif, kaitkan
% dengan Bab 2, akui keterbatasan, bahas implikasi.

\section{Implementasi Sistem}
...
\section{Hasil Efektivitas Verifikasi}
...
\section{Hasil Efisiensi}
...
\section{Hasil Reliabilitas}
...
\section{Pembahasan}
...
\section{Keterbatasan Penelitian}
...

% !TEX root = ../main.tex
\chapter{KESIMPULAN DAN SARAN}
\label{bab:kesimpulan}
% Pedoman 9.5: kesimpulan menjawab setiap rumusan masalah secara langsung,
% tanpa data baru. Saran spesifik, realistis, dan menyebut siapa pelaksananya.

\section{Kesimpulan}
\begin{enumerate}
  \item ...
\end{enumerate}

\section{Saran}
\begin{enumerate}
  \item ...
\end{enumerate}


% ---------------------------------------------------------------------------
% DAFTAR REFERENSI
% ---------------------------------------------------------------------------
\clearpage
\begingroup
\setstretch{1}
\printbibliography[heading=bibintoc,title={DAFTAR REFERENSI}]
\endgroup

% ---------------------------------------------------------------------------
% LAMPIRAN (jika ada)
% ---------------------------------------------------------------------------
\uilampiran
% !TEX root = ../main.tex
\chapter{PROMPT AGEN VERIFIKASI}
\label{lamp:prompt}
...

\chapter{CONTOH KELUARAN JSON}
\label{lamp:json}
...


\end{document}
